\providecommand{\Res}[2]{\qty[round-mode=figures,round-precision=3]{\fpeval{#1}}{#2}}
\providecommand{\Val}[1]{\num[round-mode=figures,round-precision=4]{\fpeval{#1}}}
\providecommand{\ResSci}[2]{\qty[round-mode=figures,round-precision=3,exponent-mode=scientific]{\fpeval{#1}}{#2}}
\subsubsection*{Question a) Moment magnétique}
\begin{enumerate}[label=\arabic*)]
    \item \textbf{Variables connues et recherchées :} $N = \v{N}$, $I = \qty{\v{I}}{A}$, $A = \num{\vsi{a}} \cdot \num{\vsi{b}} = \Res{\vsi{a}*\vsi{b}}{m^2}$ ; variable recherchée : $\mu$
    \item \textbf{Équations pertinentes :}
    \[ \mu = NIA \]
    \item \textbf{Isolation de la variable recherchée :} l'équation est déjà isolée.
    \item \textbf{Résolution numérique :}
    \[ \mu = \v{N} \cdot \num{\v{I}} \cdot \Val{\vsi{a}*\vsi{b}} = \Res{\v{N}*\v{I}*\vsi{a}*\vsi{b}}{A.m^2} \]
    \item \textbf{Réponse :} $\mu = \mathbf{\Res{\v{N}*\v{I}*\vsi{a}*\vsi{b}}{A.m^2}}$.
\end{enumerate}

\subsubsection*{Question b) Moment de force}
\begin{enumerate}[label=\arabic*)]
    \item \textbf{Variables connues et recherchées :} $\mu = \Res{\v{N}*\v{I}*\vsi{a}*\vsi{b}}{A.m^2}$, $B = \qty{\v{B}}{T}$, $\theta = \qty{\v{th}}{\degree}$ ; variable recherchée : $\tau$
    \item \textbf{Équations pertinentes :}
    \[ \vec{\tau} = \vec{\mu} \times \vec{B} \qquad \tau = \mu B \sin\theta \]
    \item \textbf{Isolation de la variable recherchée :} l'équation est déjà isolée.
    \item \textbf{Résolution numérique :}
    \[ \tau = \Val{\v{N}*\v{I}*\vsi{a}*\vsi{b}} \cdot \num{\v{B}} \cdot \sin(\qty{\v{th}}{\degree}) = \Res{(\v{N}*\v{I}*\vsi{a}*\vsi{b})*\v{B}*sind(\v{th})}{N.m} \]
    \item \textbf{Réponse :} $\tau = \mathbf{\Res{(\v{N}*\v{I}*\vsi{a}*\vsi{b})*\v{B}*sind(\v{th})}{N.m}}$.
\end{enumerate}

\subsubsection*{Question c) Moment de force maximal}
Le moment est maximal lorsque $\sin\theta = 1$ (normale perpendiculaire au champ) :
\[ \tau_{\max} = \mu B = \Val{\v{N}*\v{I}*\vsi{a}*\vsi{b}} \cdot \num{\v{B}} = \mathbf{\Res{(\v{N}*\v{I}*\vsi{a}*\vsi{b})*\v{B}}{N.m}} \]

\subsubsection*{Question d) Travail pour tourner la bobine}
\begin{enumerate}[label=\arabic*)]
    \item \textbf{Variables connues et recherchées :} $\mu = \Res{\v{N}*\v{I}*\vsi{a}*\vsi{b}}{A.m^2}$, $B = \qty{\v{B}}{T}$, $\theta_i = \qty{0}{\degree}$, $\theta_f = \qty{\v{th}}{\degree}$ ; variable recherchée : $W$
    \item \textbf{Équations pertinentes :}
    \[ U = -\vec{\mu}\cdot\vec{B} = -\mu B\cos\theta \qquad W_{\text{ext}} = \Delta U \]
    \item \textbf{Isolation de la variable recherchée :}
    \[ W_{\text{ext}} = -\mu B\cos\theta_f - (-\mu B\cos\theta_i) = \mu B\,(1 - \cos\theta_f) \]
    \item \textbf{Résolution numérique :}
    \[ W_{\text{ext}} = \Val{\v{N}*\v{I}*\vsi{a}*\vsi{b}} \cdot \num{\v{B}} \cdot \left(1 - \cos(\qty{\v{th}}{\degree})\right) = \Res{(\v{N}*\v{I}*\vsi{a}*\vsi{b})*\v{B}*(1-cosd(\v{th}))}{J} \]
    \item \textbf{Réponse :} $W_{\text{ext}} = \mathbf{\Res{(\v{N}*\v{I}*\vsi{a}*\vsi{b})*\v{B}*(1-cosd(\v{th}))}{J}}$.
\end{enumerate}
